\documentclass[11pt,a4paper]{article}
\usepackage[utf8]{inputenc}
\usepackage[margin=1in]{geometry}
\usepackage{amsmath,amssymb,amsfonts}
\usepackage{graphicx}
\usepackage{booktabs}
\usepackage{hyperref}
\usepackage{array}
\usepackage{xcolor}
\usepackage{microtype}
\usepackage{caption}
\usepackage{subcaption}

\hypersetup{
    colorlinks=true,
    linkcolor=blue!70!black,
    citecolor=blue!70!black,
    urlcolor=blue!70!black
}

\title{\textbf{Inventory Demand Forecasting via Local State-Space Kalman Filtering and Global Gradient Boosted Decision Trees}\\
\large Applied Machine Learning (CSAI) --- Technical Evaluation Report}

\author{\textbf{Team 5} \\
Department of Computer Science and Artificial Intelligence \\
School of Technology, Rishihood University \\
Sonipat, Haryana, India \\
\texttt{https://github.com/Jgauri24/aml-project}}

\date{\today}

\begin{document}

\maketitle

\begin{abstract}
Accurate multi-horizon inventory demand forecasting is foundational to supply-chain resilience, working-capital efficiency, and stockout prevention. Retail time-series data, however, violate standard Gaussian regression assumptions due to pervasive intermittency (69.33\% zero-demand frequency), weekly calendar cycles, and promotional demand shocks. In this report, we present a comprehensive benchmark comparing two opposing paradigms: \textbf{local state-space models} (Kalman Filtering variants) and \textbf{global machine learning} (Direct Multi-Horizon LightGBM), alongside classical seasonal heuristic baselines. Evaluated on a deterministic cohort of 240 series from the Walmart M5 dataset over a four-fold rolling-origin backtest ($H=28$ days, 26,880 out-of-sample evaluated predictions), our results reveal distinct operational strengths:
(1) Global LightGBM achieves superior point forecast accuracy (\textbf{0.9227 MAE}, \textbf{0.7640 WAPE}) by pooling cross-series features across the retail hierarchy;
(2) Kalman Local Level with Seasonal Harmonics achieves the lowest squared error (\textbf{2.0535 RMSE}), dampening volatile weekend demand spikes; and
(3) Local Kalman filters deliver closed-form, well-calibrated 80\% prediction intervals (\textbf{81.08\% empirical coverage}). We synthesize these findings into an operational hybrid framework and an interactive Streamlit diagnostic platform.
\end{abstract}

\vspace{1em}
\noindent\textbf{Keywords:} Demand Forecasting, State-Space Models, Kalman Filtering, Maximum Likelihood Estimation, Gradient Boosting, LightGBM, Rolling Backtesting, Uncertainty Quantification.

\newpage
\tableofcontents
\newpage

\section{Introduction}
Retail inventory management requires forecasting systems that address two core operational needs:
\begin{enumerate}
    \item \textbf{Expected Demand Sizing (Point Forecasting):} Accurate conditional expectation or median estimates to optimize purchase orders, logistics, and replenishment schedules.
    \item \textbf{Buffer Stock Planning (Uncertainty Quantification):} Well-calibrated prediction intervals to calculate safety stock without over-allocating working capital.
\end{enumerate}

Traditional statistical time-series approaches often rely on local models (e.g., ARIMA, Croston's method, State-Space models) estimated independently per series. While local models adapt dynamically to level shifts and handle missing observations natively, they cannot share structural patterns across items. In contrast, modern machine learning approaches (e.g., LightGBM, DeepAR) pool data across thousands of items, allowing sparse products to borrow statistical strength from high-volume peers.

This study implements an end-to-end benchmarking platform comparing local Kalman filter variants against global LightGBM on the Walmart M5 dataset.

\section{Dataset Architecture \& Cutoff-Safe EDA}

\subsection{The Walmart M5 Dataset \& Cohort Selection}
The canonical data source is the Walmart M5 dataset. To establish a tractable, reproducible, and balanced experimental cohort, we selected $N = 240$ SKU-store series via stratified sampling:
\begin{itemize}
    \item \textbf{Volume Stratification:} 80 Low-, 80 Medium-, and 80 High-volume series based on trailing-365-day history up to day $d_{1801}$.
    \item \textbf{Geographic Diversity:} Covers all 3 states (California, Texas, Wisconsin) and all 10 Walmart retail stores.
    \item \textbf{Category Diversity:} Spans all 3 merchandise categories (\texttt{FOODS}, \texttt{HOUSEHOLD}, \texttt{HOBBIES}) across 7 departments.
    \item \textbf{Zero-Preservation:} Real zero-demand and intermittent patterns are preserved without artificial thresholds.
\end{itemize}

\subsection{Cutoff-Safe Exploratory Data Analysis}
All exploratory analysis was conducted strictly on history prior to day $d_{1801}$ (432,240 series-days) to prevent future data leakage:
\begin{itemize}
    \item \textbf{Extreme Intermittency:} Overall, $69.33\%$ of daily observations are zero. In the low-volume group, $89.06\%$ of days have zero sales, with a maximum consecutive zero streak of $1,640$ days.
    \item \textbf{Syntetos-Boylan Classification:} 175 Intermittent series, 54 Lumpy series, 8 Smooth series, and 3 Erratic series. Standard percentage metrics like MAPE diverge due to zero denominators, necessitating WAPE and MAE.
    \item \textbf{7-Day Cyclical Seasonality:} Autocorrelation Function (ACF) peaks at lags 7, 14, 21, and 28 ($r_7 \approx 0.1639$). Weekend sales demonstrate a $25\%\text{--}40\%$ surge over mid-week days (peak-to-trough ratio = $1.369$).
    \item \textbf{SNAP Program Lift:} Supplemental Nutrition Assistance Program benefit disbursements in the first 10 days of each month produce distinct state-specific lifts: California $+4.8\%$, Texas $+6.0\%$, and Wisconsin $+8.0\%$.
\end{itemize}

\section{State-Space Modeling \& Kalman Filtering}

\subsection{Linear Gaussian State-Space System}
The general state-space representation decomposes observed demand $y_t$ into latent state components $\boldsymbol{\alpha}_t$:
\begin{align}
y_t &= \mathbf{Z}_t \boldsymbol{\alpha}_t + \varepsilon_t, \quad &\varepsilon_t \sim \mathcal{N}(0, R) \\
\boldsymbol{\alpha}_t &= \mathbf{T}_t \boldsymbol{\alpha}_{t-1} + \boldsymbol{\eta}_t, \quad &\boldsymbol{\eta}_t \sim \mathcal{N}(\mathbf{0}, \mathbf{Q})
\end{align}
where $R$ is observation variance and $\mathbf{Q}$ is the process noise covariance matrix.

\subsection{Evaluated Kalman Variants}
We implemented five independent per-series Kalman variants:
\begin{enumerate}
    \item \textbf{Single-SKU Baseline (\texttt{kalman\_single\_sku}):} Local-level model with scalar state $\alpha_t = [\mu_t]$, $T = [1]$, and $Z = [1]$.
    \item \textbf{Local Level (\texttt{kalman\_local\_level}):} Standard random walk level:
    \begin{equation}
    y_t = \mu_t + \varepsilon_t, \quad \mu_t = \mu_{t-1} + \eta_t
    \end{equation}
    \item \textbf{Local Level + Seasonal (\texttt{kalman\_local\_level\_seasonal}):} Augments the level state with a 6-dimensional cyclical state vector $\boldsymbol{\gamma}_t$:
    \begin{equation}
    \boldsymbol{\alpha}_t = [\mu_t, \gamma_{t,1}, \dots, \gamma_{t,6}]^T, \quad \mathbf{Z} = [1, 1, 0, 0, 0, 0, 0]
    \end{equation}
    Transition block $\mathbf{S}_{6 \times 6}$ rotates weekly harmonics while enforcing $\sum_{j=0}^6 \gamma_{t-j} = 0$.
    \item \textbf{Local Trend (\texttt{kalman\_local\_trend}):} Incorporates a dynamic slope $\beta_t$ with damping factor $\rho = 0.98$:
    \begin{equation}
    \boldsymbol{\alpha}_t = \begin{bmatrix} \mu_t \\ \beta_t \end{bmatrix}, \quad \mathbf{T} = \begin{bmatrix} 1 & 1 \\ 0 & \rho \end{bmatrix}, \quad \mathbf{Z} = [1, 0]
    \end{equation}
    \item \textbf{Local Trend + Seasonal (\texttt{kalman\_local\_trend\_seasonal}):} Full 8-dimensional structural state-space model.
\end{enumerate}

\subsection{Recursive Filtering and Missing Data Handling}
For state estimate $\mathbf{a}_{t|t-1}$ and error covariance $\mathbf{P}_{t|t-1}$:
\begin{align}
\text{Prediction:} \quad &\mathbf{a}_{t|t-1} = \mathbf{T} \mathbf{a}_{t-1|t-1}, \quad \mathbf{P}_{t|t-1} = \mathbf{T} \mathbf{P}_{t-1|t-1} \mathbf{T}^T + \mathbf{Q} \\
\text{Correction:} \quad &\mathbf{a}_{t|t} = \mathbf{a}_{t|t-1} + \mathbf{K}_t (y_t - \mathbf{Z} \mathbf{a}_{t|t-1}) \\
&\mathbf{P}_{t|t} = (\mathbf{I} - \mathbf{K}_t \mathbf{Z}) \mathbf{P}_{t|t-1}
\end{align}
where Kalman Gain $\mathbf{K}_t = \mathbf{P}_{t|t-1} \mathbf{Z}^T (\mathbf{Z} \mathbf{P}_{t|t-1} \mathbf{Z}^T + R)^{-1}$.

\paragraph{Stockout Invariant:} When sales $y_t$ are missing or censored (\texttt{NaN}), the measurement update is skipped ($\mathbf{a}_{t|t} = \mathbf{a}_{t|t-1}$ and $\mathbf{P}_{t|t} = \mathbf{P}_{t|t-1}$). The filter coasts on the transition model without artificial zero-imputation.

\subsection{Parameter Estimation \& Multi-Step Forecast}
Parameters $\boldsymbol{\theta} = (\ln Q, \ln R)$ are fitted per-series via Maximum Likelihood Estimation (MLE) using the prediction-error decomposition:
\begin{equation}
\ell(\boldsymbol{\theta}) = \frac{1}{2}\sum_{t=1}^T \left[ \ln(2\pi S_t) + \frac{v_t^2}{S_t} \right]
\end{equation}
The multi-step point forecast is invariant: $\mathbb{E}[y_{T+h} \mid \mathcal{F}_T] = \mu_{T|T}$. Predictive variance grows linearly:
\begin{equation}
V_h = \text{Var}(y_{T+h} \mid \mathcal{F}_T) = P_{T|T} + h \cdot Q + R
\end{equation}
Quantile intervals are zero-clipped: $\hat{q}_\tau(y_{T+h}) = \max(0, \hat{y}_{T+h} + z_\tau \sqrt{V_h})$.

\section{Feature Engineering \& Global LightGBM}

\subsection{Cross-Series Statistical Learning}
Unlike local models, global LightGBM pools all 240 series into a unified tabular feature space. This enables sparse, low-volume items to borrow seasonal and promotional patterns from high-volume items.

\subsection{Leakage-Safe Feature Pipeline}
We engineered 26 tabular features strictly at origin $T$ ($t \le T$):
\begin{itemize}
    \item \textbf{Lags:} Sales at $T-1, T-7, T-14, T-28$.
    \item \textbf{Rolling Statistics:} 7-day and 28-day rolling means and standard deviations.
    \item \textbf{Calendar Encodings:} Day-of-week, day-of-month, month, weekend binary flag, and $\sin/\cos$ cyclical transforms.
    \item \textbf{Exogenous Indicators:} State SNAP flags (\texttt{snap\_CA}, \texttt{snap\_TX}, \texttt{snap\_WI}), national event flags, and sell prices.
    \item \textbf{Temporal Invariant:} For horizon $h > \text{lag}$, future values are treated strictly as \texttt{NaN} and processed via LightGBM's native missingness splits. Tested via \texttt{assert\_no\_future\_leakage()}.
\end{itemize}

\subsection{LightGBM Configuration}
Trained with L1 loss (\texttt{regression\_l1} / MAE), predicting the conditional median. Tuned via rolling validation on the first 3 folds: \texttt{learning\_rate}: $0.05$, \texttt{n\_estimators}: $500$, \texttt{num\_leaves}: $31$, \texttt{feature\_fraction}: $0.8$.

\section{Experimental Protocol}
Evaluation is conducted using an expanding four-fold rolling backtest ($H = 28$ days):
\begin{itemize}
    \item \textbf{Fold 1:} Origin $d_{1801}$ (2016-01-04 to 2016-01-31, 6,720 forecasts)
    \item \textbf{Fold 2:} Origin $d_{1829}$ (2016-02-01 to 2016-02-28, 6,720 forecasts)
    \item \textbf{Fold 3:} Origin $d_{1857}$ (2016-02-29 to 2016-03-27, 6,720 forecasts)
    \item \textbf{Fold 4:} Origin $d_{1885}$ (2016-03-28 to 2016-04-24, 6,720 forecasts)
\end{itemize}
Total out-of-sample predictions: $\mathbf{26,880}$. Days $d_{1914}\text{--}d_{1941}$ remain a strictly sealed holdout.

\section{Experimental Results \& Discussion}

\begin{table}[htbp]
\centering
\caption{Pooled Performance Across All Folds (26,880 Observations)}
\label{tab:main_results}
\begin{tabular}{llccccccc}
\toprule
\textbf{Model Architecture} & \textbf{Paradigm} & \textbf{MAE $\downarrow$} & \textbf{RMSE $\downarrow$} & \textbf{WAPE $\downarrow$} & \textbf{PL Q10 $\downarrow$} & \textbf{PL Q50 $\downarrow$} & \textbf{PL Q90 $\downarrow$} & \textbf{80\% Cov.} \\
\midrule
\textbf{LightGBM (Global)} & Global GBDT & \textbf{0.9227} & 2.1081 & \textbf{0.7640} & --- & --- & --- & --- \\
\textbf{Kalman LL Seasonal} & State-Space & 1.0007 & \textbf{2.0535} & 0.8286 & 0.1318 & 0.5004 & 0.3306 & 85.72\% \\
Kalman LT Seasonal & State-Space & 1.0030 & 2.0576 & 0.8305 & 0.1347 & 0.5015 & 0.3304 & 85.43\% \\
Kalman Local Level & State-Space & 1.0106 & 2.0613 & 0.8368 & 0.1315 & 0.5053 & 0.3366 & 85.64\% \\
\textbf{Kalman Single-SKU (M6)} & State-Space & 1.0115 & 2.0776 & 0.8375 & 0.1426 & 0.5058 & 0.3518 & \textbf{81.08\%} \\
Kalman Local Trend & State-Space & 1.0225 & 2.0976 & 0.8466 & 0.1419 & 0.5112 & 0.3371 & 85.37\% \\
Naive Weekly Seasonal & Heuristic & 1.1869 & 2.5487 & 0.9828 & --- & --- & --- & --- \\
Naive Last Value & Heuristic & 1.2479 & 2.5151 & 1.0332 & --- & --- & --- & --- \\
\bottomrule
\end{tabular}
\end{table}

\subsection{Comparative Analysis}
\begin{itemize}
    \item \textbf{Point Accuracy:} LightGBM achieves the best MAE ($0.9227$) and WAPE ($0.7640$), outperforming naive seasonal heuristics by $26.1\%$ and Kalman filters by $7.8\%$.
    \item \textbf{Squared Volatility:} Kalman Local Level + Seasonal achieves the lowest RMSE ($2.0535$), outperforming LightGBM ($2.1081$). Harmonic state tracking absorbs weekend demand surges without generating extreme tail errors.
    \item \textbf{Uncertainty Calibration:} Single-SKU Kalman delivers $81.08\%$ empirical coverage for its nominal 80\% prediction interval.
    \item \textbf{Trend Degeneracy:} Adding linear trends (\texttt{kalman\_local\_trend}) worsened performance ($0.8466$ WAPE) due to overfitting on transient surges.
\end{itemize}

\subsection{Performance by Volume Band and Rolling Folds}

\begin{table}[htbp]
\centering
\caption{Kalman Local Level + Seasonal Performance by Volume Band}
\label{tab:vol_band}
\begin{tabular}{lcccc}
\toprule
\textbf{Volume Band} & \textbf{Observations} & \textbf{MAE $\downarrow$} & \textbf{RMSE $\downarrow$} & \textbf{WAPE $\downarrow$} \\
\midrule
High Volume & 8,960 & 1.9453 & 3.3110 & 0.7166 \\
Medium Volume & 8,960 & 0.7002 & 1.1062 & 1.1136 \\
Low Volume & 8,960 & 0.3566 & 0.6813 & 1.2740 \\
\bottomrule
\end{tabular}
\end{table}

\begin{table}[htbp]
\centering
\caption{MAE Progression Across Rolling Folds}
\label{tab:fold_progression}
\begin{tabular}{lcccc}
\toprule
\textbf{Model} & \textbf{Fold 1801} & \textbf{Fold 1829} & \textbf{Fold 1857} & \textbf{Fold 1885} \\
\midrule
LightGBM & 0.8821 & 0.9306 & 0.9103 & 0.9679 \\
Kalman LL Seasonal & 0.9723 & 1.0145 & 0.9903 & 1.0257 \\
Naive Weekly Seasonal & 1.1460 & 1.1835 & 1.1676 & 1.2507 \\
Naive Last Value & 1.2314 & 1.2528 & 1.2509 & 1.2564 \\
\bottomrule
\end{tabular}
\end{table}

\section{Interactive System Platform (\texttt{app.py})}
To operationalize our findings, we deployed an interactive Streamlit platform featuring:
\begin{enumerate}
    \item \textbf{Live Kalman Simulation:} Real-time slider manipulation of process noise $Q$ and observation noise $R$ with dynamic 80\% prediction interval updates.
    \item \textbf{Leakage-Safe Feature Explorer:} Interactive feature vector synthesis at origin $T$.
    \item \textbf{Benchmark Visualizer:} Live metric sorting and volume-band filtering across all models.
\end{enumerate}

\section{Conclusion \& Operational Recommendations}
We recommend a \textbf{hybrid operational policy}:
\begin{itemize}
    \item \textbf{Procurement Baseline:} Deploy Global LightGBM for baseline purchase order sizing (lowest linear error: 0.7640 WAPE).
    \item \textbf{Safety-Stock Sizing:} Use Kalman State-Space predictive variance $V_h$ for dynamic safety stock buffers (81.08\% empirical coverage).
\end{itemize}

\begin{thebibliography}{00}
\bibitem{b1} R. E. Kalman, ``A new approach to linear filtering and prediction problems,'' \textit{Journal of Basic Engineering}, vol. 82, no. 1, pp. 35--45, 1960.
\bibitem{b2} D. Salinas, V. Flunkert, J. Gasthaus, and T. Januschowski, ``DeepAR: Probabilistic forecasting with autoregressive recurrent networks,'' \textit{International Journal of Forecasting}, vol. 36, no. 3, pp. 1181--1191, 2020.
\bibitem{b3} A. C. Harvey, \textit{Forecasting, Structural Time Series Models and the Kalman Filter}. Cambridge University Press, 1990.
\bibitem{b4} G. Ke et al., ``LightGBM: A highly efficient gradient boosting decision tree,'' \textit{Advances in Neural Information Processing Systems (NeurIPS)}, vol. 30, 2017.
\bibitem{b5} M. Makridakis, E. Spiliotis, and V. Assimakopoulos, ``The M5 accuracy competition: Results, findings, and conclusions,'' \textit{International Journal of Forecasting}, vol. 38, no. 4, pp. 1346--1364, 2022.
\bibitem{b6} J. E. Boylan and A. A. Syntetos, \textit{Intermittent Demand Forecasting: Context, Methods and Applications}. John Wiley \& Sons, 2021.
\end{thebibliography}

\end{document}
